\section{Operaciones en los reales}

\begin{note}
  A cada \(q \in \Q\) se le asocia la clase \(\tilde{q}\) de la sucesión constante \((q, q, q, \dots)\). La función \(q \mapsto \tilde{q}\) es inyectiva, y de esta manera consideramos \(\Q \subset \R\).
\end{note}

\begin{definition}
  Sean \(s = [(a_n)]\) y \(t = [(b_n)]\) en \(\R\). Definimos
  \begin{itemize}
    \item \(s + t = [(a_n + b_n)]\),
    \item \(s \cdot t = [(a_n b_n)]\).
  \end{itemize}
\end{definition}

Hay que verificar que las definiciones son correctas.

\begin{itemize}
  \item \emph{Las sucesiones \((a_n + b_n)\) y \((a_n b_n)\) son de Cauchy.} Para la suma es inmediato. Para el producto, si \(|a_n|, |b_n| \leq M\) para todo \(n\) (las sucesiones de Cauchy son acotadas), entonces
        \[
          a_n b_n - a_m b_m = a_n (b_n - b_m) + b_m (a_n - a_m)
        \]
        y por lo tanto \(|a_n b_n - a_m b_m| \leq M \bigl( |b_n - b_m| + |a_n - a_m| \bigr)\), que es menor que \(\e\) para \(n, m\) grandes.
  \item \emph{No dependen de los representantes.} Sean \([(a_n)] = [(c_n)]\) y \([(b_n)] = [(d_n)]\), es decir, \(a_n - c_n \to 0\) y \(b_n - d_n \to 0\). Para la suma,
        \[
          (a_n + b_n) - (c_n + d_n) = (a_n - c_n) + (b_n - d_n) \to 0.
        \]
        Para el producto,
        \[
          a_n b_n - c_n d_n = b_n (a_n - c_n) + c_n (b_n - d_n).
        \]
        Como \((b_n)\) y \((c_n)\) son acotadas, existe \(M\) con \(|b_n|, |c_n| \leq M\) para todo \(n\), de donde \(|a_n b_n - c_n d_n| \leq M \bigl( |a_n - c_n| + |b_n - d_n| \bigr) \to 0\).
\end{itemize}

\begin{prop}
  \(\R\) con estas operaciones es un cuerpo.
\end{prop}

\begin{proof}
  Los axiomas de cuerpo (asociatividad, conmutatividad, distributividad, neutros \(\tilde{0}\) y \(\tilde{1}\), opuesto \(-[(a_n)] = [(-a_n)]\)) se heredan de \(\Q\) operando término a término. Queda como ejercicio. La existencia del inverso multiplicativo se prueba en el siguiente teorema.
\end{proof}

\begin{lemma}
  Sea \((a_n)\) una sucesión de Cauchy con \(a_n \not\to 0\). Entonces existen \(\delta \in \Q\), \(\delta > 0\), y \(n_0 \in \N\) tales que \(|a_n| > \delta\) para todo \(n > n_0\). Además, o bien \(a_n > \delta\) para todo \(n > n_0\), o bien \(a_n < -\delta\) para todo \(n > n_0\).
\end{lemma}

\begin{proof}
  Como \(a_n \not\to 0\), existe \(\e_1 > 0\) tal que \(|a_n| \geq \e_1\) para infinitos valores de \(n\). Como \((a_n)\) es de Cauchy, existe \(n_0\) tal que \(|a_n - a_m| < \frac{\e_1}{2}\) para todos \(n, m > n_0\). Elegimos \(M > n_0\) con \(|a_M| \geq \e_1\). Entonces, para \(n > n_0\),
  \[
    |a_n| \geq |a_M| - |a_n - a_M| > \e_1 - \frac{\e_1}{2} = \frac{\e_1}{2}.
  \]
  Sea \(\delta = \frac{\e_1}{2}\). Si \(a_M > 0\), entonces \(a_n > a_M - \frac{\e_1}{2} \geq \delta\) para \(n > n_0\); si \(a_M < 0\), análogamente \(a_n < a_M + \frac{\e_1}{2} \leq -\delta\).
\end{proof}

\begin{theorem}
  Si \(s \in \R\), \(s \neq 0\), entonces existe \(t \in \R\) tal que \(s \cdot t = 1\).
\end{theorem}

\begin{proof}
  Sea \(s = [(a_n)]\). Como \(s \neq \tilde{0}\), se tiene \(a_n \not\to 0\). Por el lema, existen \(\delta > 0\) y \(n_0\) tales que \(|a_n| > \delta\) para \(n > n_0\); en particular \(a_n \neq 0\) para \(n > n_0\), aunque puede haber finitos términos nulos al principio. Definimos
  \[
    b_n =
    \begin{cases}
      0              & \text{si } n \leq n_0, \\
      \dfrac{1}{a_n} & \text{si } n > n_0.
    \end{cases}
  \]
  La sucesión \((b_n)\) es de Cauchy: para \(n, m > n_0\),
  \[
    |b_n - b_m| = \frac{|a_m - a_n|}{|a_n| |a_m|} < \frac{|a_n - a_m|}{\delta^2},
  \]
  que es arbitrariamente pequeño para \(n, m\) grandes. Además \(a_n b_n = 1\) para \(n > n_0\), de modo que \((1, 1, \dots) - (a_n b_n) \to 0\) y \([(a_n b_n)] = \tilde{1}\). Es decir, \(t = [(b_n)]\) cumple \(s \cdot t = 1\).
\end{proof}

\section{Cuerpo ordenado}

Para ver que \(\R\) es un cuerpo ordenado hay que decir quiénes son los positivos.

\begin{definition}
  Un real \(s \in \R\) es \emph{positivo} si \(s \neq 0\) y \(s = [(a_n)]\) con \(a_n > 0\) para todo \(n\) mayor que cierto \(n_0\); es decir, si todos los términos son positivos a partir de cierto punto. Escribimos \(s > t\) si \(s - t\) es positivo.
\end{definition}

\begin{ejercicio}
  Probar que la definición no depende del representante elegido, que la suma y el producto de positivos son positivos, y que para todo \(s \in \R\) vale exactamente una de las opciones \(s > 0\), \(s = 0\), \(-s > 0\). (Sugerencia: usar el lema anterior.)
\end{ejercicio}

Veamos, a modo de ejemplo, cómo se prueban las propiedades de orden.

\begin{theorem}
  Si \(s, t, r \in \R\) y \(s > t\), entonces \(s + r > t + r\).
\end{theorem}

\begin{proof}
  Sean \(s = [(a_n)]\), \(t = [(b_n)]\) y \(r = [(c_n)]\). Como \((a_n + c_n) - (b_n + c_n) = a_n - b_n\) para todo \(n\),
  \[
    (s + r) - (t + r) = [(a_n - b_n)] = s - t > 0. \qedhere
  \]
\end{proof}

\begin{theorem}
  \(\R\) es arquimediano: si \(s, t > 0\), existe \(m \in \N\) tal que \(m \cdot s > t\).
\end{theorem}

\begin{proof}
  Sean \(s = [(a_n)]\) y \(t = [(b_n)]\). Como \(s > 0\), por el lema existen \(\delta \in \Q\), \(\delta > 0\), y \(n_0\) tales que \(a_n > \delta\) para \(n > n_0\). Como \((b_n)\) es de Cauchy, es acotada: existe \(B \in \Q\) con \(b_n \leq B\) para todo \(n\). Como \(\Q\) es arquimediano, existe \(m \in \N\) con \(m \delta > B + 1\). Entonces, para \(n > n_0\),
  \[
    m a_n - b_n > m \delta - B > 1.
  \]
  Luego \(m a_n - b_n \not\to 0\) y es positivo a partir de \(n_0\), es decir, \(m s - t > 0\).
\end{proof}

\begin{theorem}
  \(\Q\) es denso en \(\R\): dados \(r \in \R\) y \(\e \in \Q\), \(\e > 0\), existe \(q \in \Q\) tal que \(|r - q| < \e\).
\end{theorem}

\begin{proof}
  Sea \(r = [(a_n)]\) con \((a_n) \subset \Q\) de Cauchy. Existe \(n_0\) tal que \(|a_n - a_m| < \frac{\e}{2}\) para todos \(n, m > n_0\). Fijamos \(l > n_0\) y sea \(q = \tilde{a_l} = [(a_l, a_l, \dots)]\). Entonces
  \[
    r - q = [(a_n - a_l)_n], \qquad \e - (r - q) = [(\e - (a_n - a_l))_n],
  \]
  y para \(n > n_0\) se cumple \(|a_n - a_l| < \frac{\e}{2}\), de modo que \(\e - (a_n - a_l) > \frac{\e}{2} > 0\) y \(\e + (a_n - a_l) > \frac{\e}{2}\) (estas sucesiones no tienden a \(0\)). Luego \(r - q < \e\) y \(q - r < \e\), es decir, \(|r - q| < \e\).
\end{proof}

\section{Propiedad del supremo}

Sea \(S \subset \R\) un conjunto no vacío y acotado superiormente por \(M \in \R\). Fijamos \(s_0 \in S\) y construimos recursivamente dos sucesiones \((u_n)_{n \geq 0}\) y \((l_n)_{n \geq 0}\) con \(u_0 = M\) y \(l_0 = s_0\). Dados \(u_n\) y \(l_n\), sea \(m_n = \frac{l_n + u_n}{2}\) el punto medio.
\begin{enumerate}[label=(\roman*)]
  \item Si \(m_n\) es cota superior de \(S\), definimos \(u_{n+1} = m_n\) y \(l_{n+1} = l_n\).
  \item Si \(m_n\) no es cota superior de \(S\), definimos \(u_{n+1} = u_n\) y \(l_{n+1} = m_n\).
\end{enumerate}

Por inducción se verifica que para todo \(n\):
\begin{enumerate}[label=(\alph*)]
  \item \(u_n\) es cota superior de \(S\);
  \item existe \(s_n \in S\) con \(l_n \leq s_n\) (en el caso (ii), \(m_n\) no es cota superior, y existe \(s \in S\) con \(s > m_n\));
  \item \(l_n \leq s_n \leq u_n\), y además \(l_n \leq l_{n+1} \leq u_{n+1} \leq u_n\);
  \item \(u_n - l_n = \dfrac{M - s_0}{2^n}\), pues en ambos casos \(u_{n+1} - l_{n+1} = \frac{u_n - l_n}{2}\).
\end{enumerate}

\begin{lemma}
  Las sucesiones \((u_n)\) y \((l_n)\) son de Cauchy en \(\R\).
\end{lemma}

\begin{proof}
  Sea \(\e > 0\). Como \(\R\) es arquimediano y \(2^n \geq n\), existe \(n_0\) tal que \(\frac{M - s_0}{2^n} < \e\) para todo \(n > n_0\). Si \(m \geq n > n_0\), por (c) se tiene \(l_n \leq l_m \leq u_m \leq u_n\), de donde
  \[
    0 \leq l_m - l_n \leq u_n - l_n < \e, \qquad 0 \leq u_n - u_m \leq u_n - l_n < \e. \qedhere
  \]
\end{proof}

\begin{lemma}
  Toda sucesión de Cauchy \((u_n)\) en \(\R\) converge a algún \(u \in \R\).
\end{lemma}

\begin{proof}
  Como \(\Q\) es denso en \(\R\), para cada \(n\) existe \(q_n \in \Q\) con \(|u_n - q_n| < \frac{1}{n}\). Veamos que \((q_n)\) es de Cauchy en \(\Q\). Dado \(\e > 0\) racional, existe \(n_0\) tal que \(|u_n - u_m| < \frac{\e}{3}\) para \(n, m > n_0\), y por la arquimedianidad existe \(n_1\) con \(\frac{1}{n} < \frac{\e}{3}\) para \(n > n_1\). Si \(n, m > \max\{n_0, n_1\}\),
  \[
    |q_n - q_m| \leq |q_n - u_n| + |u_n - u_m| + |u_m - q_m| < \e.
  \]
  Sea \(u = [(q_n)] \in \R\). Para cada \(n\), \(|u_n - u| \leq |u_n - \tilde{q_n}| + |\tilde{q_n} - u|\). El primer término es menor que \(\frac{1}{n}\). Para el segundo, \(\tilde{q_n} - u = [(q_n - q_k)_k]\) y, dado \(\e > 0\), existe \(n_2\) tal que \(|q_n - q_k| < \e\) para \(n, k > n_2\), de modo que \(|\tilde{q_n} - u| \leq \e\) para \(n > n_2\). Por lo tanto \(u_n \to u\).
\end{proof}

\begin{lemma}
  Existe \(u \in \R\) tal que \(u_n \to u\) y \(l_n \to u\).
\end{lemma}

\begin{proof}
  Por los lemas anteriores existe \(u\) con \(u_n \to u\). Como \(u_n - l_n = \frac{M - s_0}{2^n} \to 0\), también \(l_n \to u\).
\end{proof}

\begin{theorem}
  \(\R\) tiene la propiedad del supremo: todo subconjunto no vacío y acotado superiormente tiene supremo.
\end{theorem}

\begin{proof}
  Probamos que el \(u\) del lema anterior es \(\sup S\).

  \emph{\(u\) es cota superior.} Si existiera \(s \in S\) con \(u < s\), sea \(\e = s - u > 0\). Como \(u_n \to u\), existe \(n\) con \(u_n - u < \e\), es decir, \(u_n < u + \e = s\), en contradicción con que \(u_n\) es cota superior de \(S\).

  \emph{\(u\) es la menor cota superior.} Sea \(c < u\). Como \(l_n \to u\), existe \(n\) con \(l_n > c\). Por (b), \(s_n \geq l_n > c\) con \(s_n \in S\), de modo que \(c\) no es cota superior de \(S\).
\end{proof}
